The interactor first prints an integer $T$, the number of test cases.

For each test case, the walk starts at the root of the tree.

To move through the tree, you may print one of the following queries:

\begin{itemize}
\item $\texttt{? L}$ --- attempt to move from the current node to its left child;
\item $\texttt{? R}$ --- attempt to move from the current node to its right child;
\item $\texttt{? B}$ --- attempt to move from the current node to its parent.
\end{itemize}

After each query, the interactor prints an integer:

\begin{itemize}
\item $1$ if the requested move is possible. In this case, the walk moves to the corresponding node;
\item $0$ if the requested move is impossible. In this case, the walk stays at the current node.
\end{itemize}

In particular, the query $\texttt{? B}$ returns $0$ when the walk is currently at the root.

You may perform at most $8000$ queries for each test case.

When you have determined the total number of nodes, print

$$\texttt{! N}$$

where $N$ is the total number of nodes in the current test case.

Printing the final answer does not count as a query.

If the answer is correct, the interaction proceeds to the next test case. Otherwise, your submission receives a Wrong Answer verdict.

After printing every query, do not forget to flush the output. You must also flush the output after printing the final answer.
